\clearpage\fancyhead[L]{\small PALLAS.}\hypertarget{native-327-lower}{}\label{native-327-lower}
\newcommand{\ObservedDeclinationMark}{\hbox{\vrule width0pt height6.625000bp depth0pt\special{pdf:bcontent}\special{pdf:code q 0 g 0.333333 6.625000 m 0.333333 3.791667 l 0.166667 3.750000 l 0.125000 3.541667 l 0.000000 3.541667 l 0.000000 3.041667 l 0.125000 3.041667 l 0.166667 2.875000 l 0.333333 2.833333 l 0.333333 0.916667 l 0.166667 0.875000 l 0.166667 0.166667 l 0.291667 0.166667 l 0.333333 0.000000 l 0.833333 0.000000 l 0.875000 0.166667 l 1.041667 0.166667 l 1.041667 0.875000 l 0.833333 0.916667 l 0.833333 2.833333 l 1.041667 2.875000 l 1.083333 3.041667 l 3.000000 3.041667 l 3.041667 2.875000 l 3.541667 2.875000 l 3.583333 3.041667 l 3.916667 3.041667 l 3.916667 3.208333 l 4.083333 3.250000 l 4.083333 3.375000 l 4.250000 3.416667 l 4.250000 3.541667 l 4.083333 3.583333 l 4.083333 3.750000 l 1.416667 3.750000 l 1.375000 3.541667 l 1.083333 3.541667 l 1.041667 3.750000 l 0.833333 3.791667 l 0.833333 6.083333 l 1.041667 6.125000 l 1.041667 6.458333 l 0.875000 6.458333 l 0.833333 6.625000 l h f* Q}\special{pdf:econtent}\kern4.250000bp}}

\begin{center}\textbf{PALLAS.}\par\vspace{4pt}\hrule\vspace{4pt}\fontsize{10}{12}\selectfont Göttingische gelehrte Anzeigen. Issues 8 and 9. pp. 73\,..\,77. 14 Jan. 1811.\par\vspace{4pt}\hrule\vspace{4pt}\end{center}

\par\indent In issue 198 of our Learned Notices last year [p. 61 of this volume], we reported on an investigation by Professor \textsc{Gauss} into the six oppositions of Pallas observed so far, which can no longer be represented by purely elliptic motion. Convinced that the increasing discrepancies were due to the perturbations of Pallas, especially those caused by Jupiter, Professor \textsc{Gauss} began calculating the perturbations due to Jupiter over the years 1803..1811 as early as last October. He wanted to see how much closer the agreement among the oppositions would be once these perturbations were applied. He also wanted to prepare for the future development of the general theory of these perturbations. This delicate and difficult work yielded a result as satisfactory as one could scarcely have expected. All six oppositions could be brought into agreement to within a few seconds in both longitude and latitude. Since Pallas was due to become visible again toward the end of 1810, it seemed worth determining in advance \SourceLeaf{328}how much the improved theory of the planet, which accounts for Jupiter’s perturbations, would alter the ephemeris for this next appearance calculated from purely elliptic elements. This was all the more interesting because Jupiter’s influence had been particularly strong at the heliocentric conjunction of 1809 and had altered the elements quite considerably. In particular, the eccentricity and mean motion had decreased sharply, the mean distance had increased, and the line of apsides had receded. The calculation showed a change in the ephemeris far greater than in any of the preceding years. Toward the end of October the right ascension was \emph{three quarters of a degree} smaller, and the southern declination \emph{six minutes} greater, than in the ephemeris printed in the October 1810 issue of the Monatl. Corresp. These differences continued to grow, as the following table shows.
\Needspace{15\baselineskip}
\begin{center}Correction to the ephemeris of Pallas for the perturbations:\end{center}
\begin{center}\fontsize{10}{12}\selectfont\setlength{\tabcolsep}{4pt}\renewcommand{\arraystretch}{1.15}\begin{tabular}{llrr|r|r}
 &  &  &  & \shortstack{Right\\ascension} & Declination\\\hline
1810. & Oct. & 26 & $0^{\text{u}}$ & $-45'$ & $+6'$\\
 & Nov. & 7 & 12 & $-52$ & $+7$\\
 & --- & 20 & 0 & $-59$ & $+9$\\
 & Dec. & 2 & 12 & $-67$ & $+11$\\
 & --- & 15 & 0 & $-75$ & $+12$\\
 & --- & 27 & 12 & $-84$ & \ObservedDeclinationMark\,$13$\\
1811 & Jan. & 9 & 0 & $-94$ & $+14$\\
 & --- & 21 & 12 & $-103$ & $+14$\\
 & Feb. & 3 & 0 & $-112$ & $+14$\\
\end{tabular}\end{center}

\par\indent So large a difference was of the utmost importance for finding Pallas, and Professor \textsc{Gauss} immediately communicated these results to several astronomers elsewhere, since the unusually unfavorable weather here prevented a search. To be quite sure in this important matter, however, Professor \textsc{Gauss} thought it best to repeat the entire calculation of the perturbations. Errors could easily have slipped in here and there in so lengthy a calculation, and the second calculation could be expected to yield still more reliable results simply because it could already use the improved theory as its basis. As it turned out, the differences were insignificant, amounting to only a few seconds. The agreement among the six oppositions remained just as good, and the necessary change to the ephemeris, apart from small details, remained just as large. To make it easier to judge the present agreement among the six oppositions and compare it with the large discrepancies obtained from purely elliptic elements (see our Notices, 1810, p. 1971) [p. 63 of this volume], we give the comparison here.
\Needspace{10\baselineskip}
\begin{center}\fontsize{10}{12}\selectfont\setlength{\tabcolsep}{4pt}\renewcommand{\arraystretch}{1.15}\begin{tabular}{r|r|r}
 & \multicolumn{2}{c}{Differences in}\\
 & \shortstack{mean\\longitude} & \shortstack{heliocentr.\\latitude}\\\hline
1803 & $+7''3$ & $+1''6$\\
1804 & $-6.8$ & $-4.0$\\
1805 & $+1.5$ & $-4.1$\\
1807 & $-3.3$ & $-3.8$\\
1808 & $-3.9$ & $-16.9$\\
1809 & $+5.0$ & $+3.5$\\
\end{tabular}\end{center}

\SourceLeaf{329}\par\indent As is well known, the latitude for 1808 was poorly observed. This warrants the conclusion that the perturbations of the motion of Pallas caused by the other planets, especially Saturn and Mars, are quite insignificant compared with those caused by Jupiter and may still be disregarded altogether for the time being.

The weather was also unfavorable to the astronomers elsewhere who had received these results. Only Mr. von \textsc{Lindenau} observed a small star with the transit instrument at the Seeberg observatory on 19 Dec.
\[\begin{gathered}16^{\text{u}}\;9^{\text{m}}\;7^{\text{s}}1\ \text{mean time}\quad\text{Right asc. }150^\circ\;20'\;12''45\\\text{estimated declination }20^\circ\;58'\text{ S}\end{gathered}\]
{\rightskip=0pt plus 1em \noindent This agreed to the minute with the position where Pallas had to be according to the corrected ephemeris. This single observation, however, which had yet to be confirmed, still gave no complete certainty that the star really was Pallas.\par}

\par\indent At this observatory, the night of 4 to 5 January was the first in nearly two months on which a search for Pallas was possible. As Professor \textsc{Gauss} could use only movable instruments, the moonlight and severe cold (11 degrees) made the search extremely difficult. Even so, Professor \textsc{Gauss} was rewarded for his efforts. He found Pallas again where he expected it from his corrected ephemeris and also managed to make a couple of comparisons with stars from the Histoire céleste. From these he derived the following position of the planet, subject to a more precise reduction later.

\begin{center}\fontsize{10}{12}\selectfont\setlength{\tabcolsep}{4pt}\renewcommand{\arraystretch}{1.15}\begin{tabular}{c|c|c}
1811 Mean time & Right ascension & Southern declination\\\hline
Jan. 4 $15^{\text{u}}\;4^{\text{m}}\;42^{\text{s}}$ & $151^\circ\;40'\;5''$ & $22^\circ\;22'\;36''$\\
\end{tabular}\end{center}

\par\indent The planet, incidentally, was still only of the 8th to 9th magnitude.

\Needspace{17\baselineskip}{\rightskip=0pt plus 1em \par\indent Astronomers elsewhere will easily find the planet using the correction to the ephemeris just given. Anyone wishing to extend the ephemeris or calculate it more precisely may use the following elements, which represent the planet’s present motion.\par}

\begin{center}\fontsize{10}{12}\selectfont\setlength{\tabcolsep}{4pt}\renewcommand{\arraystretch}{1.15}\begin{tabular}{llr}
\shortstack[l]{Epoch of mean longitude\\on the Göttingen meridian 1811} & .\;. & $126^\circ\;33'\;6''$\\
Longitude of perihelion & .\;.\;.\;.\;.\;.\;.\;.\;.\;.\;.\;.\;.\;.\;.\;. & $120\;55\;8$\\
Longitude of the ascending node & .\;.\;.\;.\;.\;.\;.\;.\;.\;.\;.\;. & $172\;32\;40$\\
Inclination of the orbit & .\;.\;.\;.\;.\;.\;.\;.\;.\;.\;.\;.\;.\;.\;.\;.\;. & $34\;35\;14$\\
Mean daily tropical motion & .\;.\;.\;.\;.\;.\;.\;,\;.\;. & $768''980$\\
Eccentricity & .\;.\;.\;.\;.\;.\;.\;.\;.\;.\;.\;.\;.\;.\;.\;.\;.\;.\;.\;. & $0.24160$\\
Logarithm of the semimajor axis & .\;.\;.\;.\;.\;.\;.\;.\;.\;.\;.\;. & $0.44278$\\
\end{tabular}\end{center}

\par\indent Professor \textsc{Gauss} believes that publishing these results without delay will be of service to practical astronomers who wish to observe Pallas. He reserves a full account of his work on the perturbations of Pallas for another occasion.
